% from ega2-1-fr.tex
\oldpage[II]{97}

\subsection{Le critère de Serre.}
\label{subsection:II.5.2-fr}

\begin{theorem}[5.2.1]
\label{II.5.2.1-fr}
\emph{(Critère de Serre.)} Soit $X$ un schéma quasi-compact ou un préschéma
dont l'espace sous-jacent est noethérien. Les conditions suivantes sont
équivalentes~:
\begin{enumerate}
  \item[\rm{a)}] $X$ est un schéma affine.
  \item[\rm{b)}] Il existe une famille d'éléments
  $f_\alpha\in A=\Gamma(X,\mathcal{O}_X)$ tels que les $X_{f_\alpha}$ soient
  affines et que l'idéal engendré par les $f_\alpha$ dans $A$ soit égal à
  $A$.
  \item[\rm{c)}] Le foncteur $\Gamma(X,\mathcal{F})$ est exact en
  $\mathcal{F}$ dans la catégorie des $\mathcal{O}_X$-Modules
  quasi-cohérents~; autrement dit, si
  \[
    0\to\mathcal{F}'\to\mathcal{F}\to\mathcal{F}''\to0 \tag{*}
  \]
  est une suite exacte de $\mathcal{O}_X$-Modules quasi-cohérents, la suite
  \[
    0\to\Gamma(X,\mathcal{F}')\to\Gamma(X,\mathcal{F})
      \to\Gamma(X,\mathcal{F}'')\to0
  \]
  est exacte.
  \item[\rm{c')}] La propriété c) a lieu pour toute suite exacte (*) de
  $\mathcal{O}_X$-Modules quasi-cohérents telle que $\mathcal{F}$ soit
  isomorphe à un sous-$\mathcal{O}_X$-Module d'un produit fini
  $\mathcal{O}_X^n$.
  \item[\rm{d)}] $\mathrm{H}^1(X,\mathcal{F})=0$ pour tout
  $\mathcal{O}_X$-Module quasi-cohérent $\mathcal{F}$.
  \item[\rm{d')}] $\mathrm{H}^1(X,\mathcal{J})=0$ pour tout faisceau
  quasi-cohérent d'idéaux $\mathcal{J}$ de $\mathcal{O}_X$.
\end{enumerate}
\end{theorem}

Il est évident que a) implique b)~; b) implique d'autre part que les
$X_{f_\alpha}$ recouvrent $X$, puisque par hypothèse la section $1$ est
combinaison linéaire des $f_\alpha$ et que les $D(f_\alpha)$ recouvrent
alors $\operatorname{Spec}(A)$. La dernière assertion de
(\hyperref[II.4.5.2-fr]{4.5.2}) implique donc que
$X\to\operatorname{Spec}(A)$ est un isomorphisme.

On sait que a) implique c) (\hyperref[I.1.3.11-fr]{I, 1.3.11}) et c)
entraîne trivialement c'). Prouvons que c') implique b). Tout d'abord, c')
entraîne que, pour tout point fermé $x\in X$ et tout voisinage ouvert $U$ de
$x$, il existe $f\in A$ tel que $x\in X_f\subset X-U$. Soit en effet
$\mathcal{J}$ (resp. $\mathcal{J}'$) le faisceau quasi-cohérent d'idéaux de
$\mathcal{O}_X$ définissant le sous-préschéma fermé réduit de $X$ ayant pour
espace sous-jacent $X-U$ (resp. $(X-U)\cup\{x\}$)
(\hyperref[I.5.2.1-fr]{I, 5.2.1})~; il est clair que l'on a
$\mathcal{J}'\subset\mathcal{J}$ et que
$\mathcal{J}''=\mathcal{J}/\mathcal{J}'$ est un $\mathcal{O}_X$-Module
quasi-cohérent ayant pour support $\{x\}$ et tel que
$\mathcal{J}''_x=k(x)$. L'hypothèse c') appliquée à la suite exacte
$0\to\mathcal{J}'\to\mathcal{J}\to\mathcal{J}''\to0$ montre que
$\Gamma(X,\mathcal{J})\to\Gamma(X,\mathcal{J}'')$ est surjectif. La section
de $\mathcal{J}''$ dont le germe en $x$ est $1_x$ est donc l'image d'une
section $f\in\Gamma(X,\mathcal{J})\subset\Gamma(X,\mathcal{O}_X)$, et on a
par définition $f(x)=1_x$ et $f(y)=0$ dans $X-U$, ce qui établit notre
assertion. D'ailleurs, si $U$ est affine, il en est de même de $X_f$
(\hyperref[I.1.3.6-fr]{I, 1.3.6}), donc la réunion des $X_f$ qui sont affines
($f\in A$) est un ensemble ouvert $Z$ contenant tous les points fermés de
$X$~; comme $X$ est un espace de Kolmogoroff quasi-compact, on a
nécessairement $Z=X$ (\hyperref[0.2.1.3-fr]{0, 2.1.3}). Puisque $X$ est
quasi-compact, il y a un nombre fini d'éléments $f_i\in A$
($1\leq i\leq n$) tels que les $X_{f_i}$ soient affines et recouvrent $X$.
Considérons alors l'homomorphisme $\mathcal{O}_X^n\to\mathcal{O}_X$ défini
par les sections $f_i$ (\hyperref[0.5.1.1-fr]{0, 5.1.1})~; comme, pour tout
$x\in X$, un au moins des $(f_i)_x$ est inversible, cet homomorphisme est
surjectif, et on a donc une suite exacte
$0\to\mathcal{R}\to\mathcal{O}_X^n\to\mathcal{O}_X\to0$, où $\mathcal{R}$
est un sous-$\mathcal{O}_X$-Module quasi-cohérent de $\mathcal{O}_X^n$. Il
résulte alors

